% !TeX program = xelatex
% UTF-8; transcription of all fourteen handwritten images.
% Medical statements are preserved as written, not revised.
\ifdefined\settextfont
  \def\OUDStandalone{0}
\else
  \def\OUDStandalone{1}
  \documentclass[11pt,a4paper]{article}
  \usepackage[margin=2cm]{geometry}
  \usepackage{amsmath,amssymb,enumitem}
  \usepackage{xepersian}
  \IfFontExistsTF{Amiri}{\settextfont{Amiri}}{\settextfont{FreeSerif}}
  \setlatintextfont{FreeSerif}
  \setlist[itemize]{itemsep=3pt,topsep=4pt}
\fi
\providecommand{\en}[1]{\lr{#1}}
\newcommand{\oudpage}[2]{\clearpage\section*{صفحهٔ #1}\noindent{\footnotesize موضوع: #2\quad تاریخ: \en{/\quad/}\hfill\en{pic#1.jpg}}\par\medskip}
\newcommand{\oudarr}{\ensuremath{\leftarrow}}
\newcommand{\oudunclear}[1]{\textit{[خوانش نامطمئن: #1]}}
\ifnum\OUDStandalone=1 \begin{document}\fi

% pic01.jpg = 20261013_024245341.jpg; handwritten page 1.
\oudpage{01}{کاپلان ـ ۱۱٫۹ ـ اختلالات مرتبط با اوپیوئید}
\section*{\en{OUD (Opioid Use Disorder)}}
۱٫۵\% مردم جهان اعتیاد به اوپیوم دارند.

\en{OUD} یک اختلال مزمن و ریلاپس‌کننده است.

در \en{DSM-5}، به ۵ تشخیص مختلف تقسیم می‌شوند:
\begin{itemize}
\item \en{OUD}.
\item مسمومیت با اوپیوئید.
\item ویتدراوال اوپیوئید.
\item \en{other opioid induced disorders}.
\item \en{unspecified opioid related disorders}.
\end{itemize}
تشخیص \en{OUD} \oudarr{} حداقل ۲ تا از ۱۱ تا کرایتریا را داشته باشد طی ۱۲ ماه اخیر.

تقسیم‌بندی کرایتریاها:
\begin{itemize}
\item \en{Impaired Control}.
\item \en{Social impairment}.
\item \en{Risky use}.
\item \en{pharmacologic changes}.
\end{itemize}
در \en{OUD} ما مصرف طولانی‌مدت و جبری (\en{Compulsive}) داریم.

% pic02.jpg = 20261013_024254101.jpg; handwritten page 2.
\oudpage{02}{کاپلان ـ ۱۱٫۹ ـ اختلالات مرتبط با اوپیوئید}
\subsection*{سایر دسته‌بندی‌های \en{OUD}}
\begin{itemize}
\item براساس مدت ترک (ریمیشن):
\begin{itemize}
\item \en{Early remission}: هیچ‌کدام از کرایتریاها وجود نداشته، بین ۳ ماه تا ۱۲ ماه.
\item \en{Sustained rem}: بیشتر از ۱۲ ماه.
\end{itemize}
\item براساس وجود درمان نگهدارنده (\en{Maintenance}): آیا بیمار داروی خاصی برای درمان \en{OUD} به صورت نگهدارنده می‌گیرد؟
\item وجود محیط کنترل‌شده: آیا بیمار در یک محیط کنترل‌شده قرار گرفته یا خیر؟
\end{itemize}
\subsection*{شدت \en{OUD}}
\begin{itemize}
\item \en{Mild}: ۲–۳ علامت.
\item \en{Moderate}: ۴–۵ علامت.
\item \en{Severe}: ۶–۱۱ علامت.
\end{itemize}
\subsection*{مسمومیت با اوپیوئید}
بیماری که اخیراً اوپیوئید مصرف کرده، علائم بالینی، تغییرات رفتاری یا سایکولوژیک پیدا می‌کند:
\begin{itemize}
\item \en{Dysphoria}.
\item مختل شدن قضاوت.
\item تغییرات سایکوموتور (آژیتاسیون / ریتاردیشن).
\item حین یا اندکی بعد از مصرف رخ می‌دهد.
\end{itemize}

% pic03.jpg = 20261013_024257961.jpg; handwritten page 3.
\oudpage{03}{کاپلان ـ ۱۱٫۹ ـ اختلالات مرتبط با اوپیوئید}
سایر یافته‌های مسمومیت:
\begin{latin}
Drowsiness or Coma, Slurred Speech, Impaired attention, Imp Memory
\end{latin}
\subsection*{ویتدراوال اوپیوئید}
بیمار پس از یک دوره مصرف اوپیوئید، مصرف خود را کاهش می‌دهد یا قطع می‌کند یا اینکه یک آنتاگونیست اوپیوئید دریافت می‌نماید.

تشخیص ویتدراوال:
\begin{itemize}
\item وجود ۳ تا از ۹ علامت ویتدراوال.
\item از نظر کلینیکی باید اهمیت داشته باشد.
\end{itemize}
اختلال مرتبط با اوپیوئیدی که هیچ‌کدام از کرایتریاها را پر نمی‌کند:
\begin{latin}
Unspecified Opioid-related disorder
\end{latin}
نباید اختلالات زیر باشد:
\begin{latin}
Substance/Medication-induced (Depressive, psychotic, anxiety, bipolar)
\end{latin}
اکثر بیماران مرد هستند (۹۰\%).

\subsection*{علائم بالینی انتاکسیکیشن اوپیوئید}
یوفوریا اولیه، آپاتی، دیس‌فوریا، تغییر سایکوموتور و مختل شدن قضاوت جایگزین می‌شود.

% pic04.jpg = 20261013_024305399.jpg; handwritten page 4.
\oudpage{04}{کاپلان ـ ۱۱٫۹ ـ اختلالات مرتبط با اوپیوئید}
\begin{itemize}
\item میوز مردمک‌ها، دروزینس و اسلورد اسپیچ.
\item سیانوز، سرد شدن پوست، افت \en{T}.
\item ادم ریه: ممکن است در \en{CXR} دیده شود.
\item \en{frothing at the mouth}.
\item کاهش \en{BP}؛ کاهش شدید \en{BP} تنها در آنوکسی دیده می‌شود.
\item آریتمی قلبی به ندرت (به دلیل آنوکسی و \en{quinine}).
\item مردمک‌ها به نور \en{reaction} می‌دهند.
\item نارسایی تنفسی و آپنه.
\end{itemize}
در مسمومیت با فنتانیل ممکن است لارنگواسپاسم رخ دهد.
\begin{latin}
Wooden Chest Syndrome
\end{latin}
\subsection*{\en{Opioid Tolerance}}
به دنبال مصرف سنگین رخ می‌دهد.

با اینکه بیمار تحمل پیدا می‌کند ولی با هر میزان مصرف حتی به میزان کمی یوفوریا پیدا می‌کند. \oudunclear{عبارت «با هر میزان مصرف حتی»}

\subsection*{کدام داروها می‌توانند باعث ویتدراوال اوپیوئید شوند؟}
\begin{itemize}
\item پارشیال آگونیست اوپیوئید: بوپرنورفین.
\item آنتاگونیست اوپیوئید: نالوکسون، نالترکسون.
\end{itemize}

% pic05.jpg = 20261013_024310197.jpg; handwritten page 5.
\oudpage{05}{\rule{4cm}{0.2pt}}
\subsection*{علائم ویتدراوال}
اضطراب، تحریک‌پذیری، \en{restlessness}، احساس درد در کمر و پا (\en{Achy pain})، افزایش حساسیت به درد.

ویتدراوال شدید \oudarr{} تهوع و استفراغ، کرامپ شکمی و اسهال، درد شدید عضلانی، اشک‌ریزش، آبریزش بینی، میدریاز، افزایش تعریق، تب و لرز، بی‌خوابی، گیجی و دلیریوم.

\subsection*{ویتدراوال \en{Spontaneous}}
\begin{itemize}
\item ۸–۹ ساعت پس از اوپیوئید کوتاه‌اثر (هروئین، مورفین، اکسی‌کدون، هیدروکدون).
\item ۲۴ ساعت پس از اوپیوئید طولانی‌اثر (متادون، بوپرنورفین، فنتانیل).
\item ظرف ۳۶–۴۸ ساعت به حداکثر شدت می‌رسد.
\item پس از ۵–۷ روز برای کوتاه‌اثر و ۲ هفته برای طولانی‌اثر رفع می‌شود.
\end{itemize}
\subsection*{\en{Precipitated}}
(دادن آنتاگونیست یا پارشیال آگونیست)
\begin{itemize}
\item ظرف ۲–۳ ساعت پس از نالترکسون و سایر فرم‌های \en{PO} یا \en{IM}.
\item ظرف چند دقیقه پس از فرم وریدی (نالوکسان).
\end{itemize}
ویتدراوال ناشی از نالترکسون از فرم \en{Spontaneous} شدیدتر است ولی سریع‌تر رفع می‌شود؛ معمولاً ظرف ۱ روز.

% pic06.jpg = 20261013_024322069.jpg; handwritten page 6.
\oudpage{06}{\rule{4cm}{0.2pt}}
\subsection*{علائم و نشانه‌های ویتدراوال موجود در کرایتریا}
\begin{enumerate}
\item مود دیسفوریک.
\item تهوع و استفراغ.
\item درد عضلانی.
\item آبریزش و اشک‌ریزش.
\item میدریاز یا پیلوارکشن یا تعریق.
\item اسهال.
\item خمیازه کشیدن.
\item تب.
\item بی‌خوابی.
\end{enumerate}
برای ویتدراوال خودبه‌خودی باید به صورت مزمن مصرف کرده باشد، ولی در فرم \en{precipitated} تنها چند روز مصرف کافی است.

از معیار \en{COWS} برای اندازه‌گیری شدت ویتدراوال اوپیوئید استفاده می‌شود.
\begin{latin}
Resting Pulse Rate, Sweating, Restlessness, pupil size, pain,\par
GI Symptoms, Tremor, Yawning, Anxiety, Piloerection
\end{latin}
یادداشت حاشیه: رینوره +.

حتی قطع ناگهانی بوپرنورفین نیز می‌تواند ویتدراوال دهد ولی معمولاً خفیف است:
\begin{itemize}
\item بی‌خوابی، تحریک‌پذیری، اضطراب، دیپرشن.
\item ظرف ۱۰–۲۱ روز رفع می‌شوند. \oudunclear{عدد انتهای بازه: ۲۱}
\item ویتدراوال ۹۶ ساعت پس از قطع شروع می‌شود چون نیمه‌عمر بالایی است.
\item به همین دلیل می‌توان به صورت اینترمیتنت استفاده کرد.
\end{itemize}

% pic07.jpg = 20261013_024330662.jpg; handwritten page 7.
\oudpage{07}{\rule{4cm}{0.2pt}}
\subsection*{سندرم ویتدراوال پست اکوت یا \en{protracted}}
\begin{itemize}
\item پس از گذشت ویتدراوال اولیه، تغییرات فیزیولوژیکی غیرطبیعی که تا هفته‌ها و ماه‌ها باقی می‌مانند.
\item مود دیپرس، انهدونیا، اضطراب، تحریک‌پذیری، مود لبیل، عدم تحمل استرس، بی‌خوابی، کراوینگ مواد.
\item به علت تغییرات در سیستم‌های سروتونرژیک و دوپامینرژیک است.
\item یکی از مهم‌ترین علل اعتیاد مجدد است.
\end{itemize}
\begin{latin}
COWS $\rightarrow$ Mild (5--12), Moderate (13--24), Modestly Severe (25--36),\par
Severe withdrawal (37--)
\end{latin}
\subsection*{درمان \en{OUD}}
\begin{itemize}
\item بهترین نگاه سه‌جانبه دارد: دارو برای نرمال نمودن عملکرد مغز.
\item مداخلات سایکوسوشال برای احیای رفتارها و ها‌بیت‌ها.
\item حمایت‌های مبتنی بر جامعه.
\item اهداف اولیهٔ درمان:
\begin{itemize}
\item کاهش یا قطع مصرف.
\item جلوگیری از عود و \en{relapse}.
\item کاهش علائم پزشکی، روان‌پزشکی و اجتماعی.
\item ارتقاء عملکرد پزشکی، روانی، اجتماعی بیمار و کیفیت زندگی.
\end{itemize}
\end{itemize}

% pic08.jpg = 20261013_024337999.jpg; handwritten page 8.
\oudpage{08}{\rule{4cm}{0.2pt}}
\subsection*{درمان \en{OUD}}
\begin{itemize}
\item هدف درمان براساس بیمار و درمانگر متفاوت است.
\item معمولاً قطع کامل مصرف در ابتدای درمان یک هدف شدنی نیست.
\item مداخلات سایکوسوشال:
\begin{itemize}
\item \en{CBT} بسیار کمک‌کننده است.
\item هرچه بیماری‌های همراه روانی بیشتر باشد \en{CBT} بهتر است.
\end{itemize}
\item مداخلات دارویی:
\begin{itemize}
\item داروهای تأییدشده: متادون، نالترکسون خوراکی، بوپرنورفین، نالترکسون طولانی‌اثر تزریقی.
\item ۱۰–۳۰\% بیمارانی که تحت درمان با متادون و \en{b2} در کلینیک سرپایی قرار می‌گیرند ظرف ۱۲ ماه ترک بدون \en{relapse} دارند.
\item بیماری که اوپیوئیدهای دارویی را سوءمصرف داشته، بهتر است با درمان \en{maintenance} بهبود درمان شود.
\end{itemize}
\end{itemize}
\subsection*{متادون}
\begin{itemize}
\item با دوزی می‌دهیم که \en{Craving} نداشته باشد: \en{80--120 mg/d}.
\item دوزهای پایین‌تر علائم ترک را رفع می‌کنند ولی \en{Craving} را نه.
\item دوز شروع \en{20--30 mg} روزانه \oudarr{} افزایش \en{10 mg} هر هفته.
\item \en{LFT} چک کنیم.
\item \en{QT} را طولانی می‌کند، وابسته به دوز.
\item در دوز بالا: نارسایی تنفسی.
\end{itemize}

% pic09.jpg = 20261013_024343093.jpg; handwritten page 9.
\oudpage{09}{\rule{4cm}{0.2pt}}
\begin{itemize}
\item سطح تستوسترون و \en{FSH} سرم را $\downarrow$.
\item بی‌خوابی و کابوس.
\item داروهای مؤثر بر \en{CYP3A4}:
\begin{itemize}
\item مثل ریفامپین، فنی‌توئین، کاربامازپین، داروناویر، افاویرنز \oudarr{} سطح متادون را $\downarrow$.
\item فلووکسامین، اریترومایسین، دلاویردین، کتوکونازول، فلوکونازول \oudarr{} سطح متادون را $\uparrow$.
\end{itemize}
\item مترونیدازول \en{QT} طولانی.
\end{itemize}
\subsection*{بوپرنورفین}
\begin{itemize}
\item در دوزهای بالا اثر آنتاگونیستی بیشتر می‌شود (\en{16 mg <}، مطابق جهت علامت در تصویر).
\item زیرزبانی باید استفاده شود؛ جذب خوراکی ضعیف است.
\item در ترکیب با نالوکسان وجود دارد \oudarr{} احتمال سوءمصرف را کم می‌کند.
\item نسبت ۴ به ۱.
\item اگر زیرزبانی مصرف شود نالوکسان جذب نمی‌شود ولی اگر تزریق شود می‌تواند \en{withdrawal} بدهد.
\item نیمه‌عمر ۲۴–۴۸ ساعت \oudarr{} روزی یک بار مصرف می‌شود.
\item \en{12 mg} بوپرنورفین معادل \en{50--60 mg} متادون اثر دارد.
\item در دوز بیشتر از \en{16 mg} اثر کمی دارد؛ نباید دوز بالاتر بدهیم.
\item باید به بیماری که علائم \en{Mild withd} دارد بدهیم؛ اگر بدهیم بدون علامت باشد علائم \en{withd} را پیدا می‌کند.
\item دوز تدریجی بالا می‌رود.
\item احتمال \en{overdose} برخلاف متادون خیلی کم است.
\end{itemize}

% pic10.jpg = 20261013_024349981.jpg; handwritten page 10.
\oudpage{10}{\rule{4cm}{0.2pt}}
\begin{itemize}
\item ترکیب \en{b2} با بنزودیازپین‌ها و سایر \en{Sedative}ها بسیار خطرناک است.
\item عوارض \en{b2} اکثراً خفیف هستند و به مرور بهتر می‌شوند.
\item تداخل دارویی مثل متادون.
\item \en{QT} را طولانی نمی‌کند.
\item اختلال شناختی ایجاد نمی‌کند.
\item دوز شروع: \en{12 mg} روزانه؛ \en{2 mg} افزایش؛ حداکثر ۲۴ (عدد قبلی خط خورده است).
\end{itemize}
مصرف اوپیوئید در بارداری برای مادر و جنین مضر است ولی تراتوژن نیست.
\begin{itemize}
\item ریسک عفونت، پره‌اکلامپسی، سقط، زایمان زودرس، \en{PROM} $\uparrow$.
\item باید مادر درمان دارویی مین‌تننس متادون یا \en{b2} باشد.
\end{itemize}
بیماری که به متادون، \en{b2} و نالترکسون جواب نمی‌دهد یا تحمل نمی‌کند؟
\begin{itemize}
\item مورفین خوراکی (آهسته‌رهش، روزی ۱ دوز).
\end{itemize}
\subsection*{به کدام بیماران نالترکسون می‌دهیم؟}
\begin{itemize}
\item عدم پاسخ به متادون و \en{b2}.
\item بیمار با انگیزه‌ای که نمی‌خواهد هیچ آگونیستی مصرف کند.
\item به درمان با آگونیست جواب داده‌اند ولی الان می‌خواهند آن را نیز قطع کنند.
\item دوز تزریقی \en{Ext rel}: \en{380 mg/IM} هر ۴ هفته.
\end{itemize}

% pic11.jpg = 20261013_024354743.jpg; handwritten page 11.
\oudpage{11}{\rule{4cm}{0.2pt}}
\begin{itemize}
\item قرص \en{50 mg} روزانه یا ۳ بار در هفته.
\item باید حداقل یک هفته از قطع اوپیوئید گذشته باشد.
\item قبل از شروع تست منفی ادرار برای اوپیوئیدها لازم است.
\item می‌توان آنتاگونیست چلنج هم انجام داد:
\begin{latin}
Amp Naloxone 0.8--1.2 mg IM\par
Tab Naltrexone 6--12.5 mg PO
\end{latin}
\item تداخل دارویی خاصی ندارد.
\item اگر بیمار دوباره اوپیوئیدی مصرف کند که \oudunclear{عبارت دست‌نویس: آنتاگونیست / ویتدراوال} می‌باشد.
\end{itemize}
\subsection*{درمان اوردوز اوپیوئید}
بیمار معمولاً بی‌هوش و دچار نارسایی تنفسی است.
\begin{itemize}
\item حفاظت از راه هوایی.
\item حمایت تنفسی.
\item اکسیژن‌درمانی با رزرو بگ.
\item نالوکسان \en{ASAP} به صورت \en{IV} یا \en{IM}.
\end{itemize}
\subsection*{دوز نالوکسان}
\begin{itemize}
\item \en{0.2--0.4 mg} در بیماری که تنفس دارد.
\item \en{1--2 mg} در بیمار آپنه یا ایست قلبی‌تنفسی.
\item نیمه‌عمر نالوکسان ۶۰–۹۰ دقیقه است.
\end{itemize}

% pic12.jpg = 20261013_024403202.jpg; handwritten page 12.
\oudpage{12}{\rule{4cm}{0.2pt}}
\begin{itemize}
\item علائم بهبودی \oudarr{} افزایش \en{RR} و دیلاتاسیون مردمک‌ها.
\item باید ظرف ۲–۳ دقیقه رخ دهد (در تزریق \en{IV}).
\item اگر جواب نداد \en{1--2 mg} می‌دهیم و هر ۲–۳ دقیقه تکرار.
\end{itemize}
اگر بیمار انتوبه‌شده و تحت \en{MV} بود دیگر نیاز به آنتی‌دوت ندارد.

هدف دادن نالوکسان \oudarr{} بازیابی \en{RR > 10} و \en{SpO2 > 92\%}.
\begin{itemize}
\item هدف، هوشیار شدن کامل نیست.
\end{itemize}
احتمال \en{ARDS} در افرادی که ناگهان آنتی‌دوت زیاد دریافت کرده‌اند، خیلی بیشتر.

اگر بعد از \en{10 mg} نالوکسان بهبودی حاصل نشد، به تشخیص دیگری فکر کنیم:
\begin{itemize}
\item هیپوگلیسمی و \en{BNZ}.
\end{itemize}
\subsection*{درمان ویتدراوال اوپیوئید (سم‌زدایی)}
اهداف سم‌زدایی:
\begin{itemize}
\item رفع علل وابستگی به اوپیوئید و جلوگیری از علائم تهدیدکنندهٔ حیات.
\item پایدارسازی سلامت جسمانی و روانی بیمار.
\item سایکواجوکیشن بیمار.
\end{itemize}

% pic13.jpg = 20261013_024408528.jpg; handwritten page 13.
\oudpage{13}{\rule{4cm}{0.2pt}}
اگر سم‌زدایی همراه با یک برنامهٔ فعال جلوگیری از ریلاپس نباشد، به احتمال ۹۰\% شکست می‌خورد.

\subsection*{اندیکاسیون بستری بیمار جهت سم‌زدایی}
\begin{itemize}
\item سابقهٔ قبلی ویتدراوال شدید یا عارضه‌دار.
\item بیماری که انتظار داشته باشیم ویتدراوال شدیدی دارد.
\item بیماری که کوموربیدیتی روان‌پزشکی و طبی دارد.
\end{itemize}
\subsection*{استراتژی‌های سم‌زدایی دارویی}
\begin{itemize}
\item تنها درمان علامتی (اوپیوئید نمی‌دهیم).
\item سم‌زدایی با یک آگونیست.
\item سم‌زدایی با یک آنتاگونیست + درمان علامتی.
\end{itemize}
\subsection*{درمان علامتی}
\begin{itemize}
\item $\alpha_2$ آدرنرژیک آگونیست‌ها: کاهش برانگیختگی سمپاتیک \oudarr{} کلونیدین.
\item دوز شروع: \en{0.1--0.3 mg TDS or qid PO}.
\item در بیماری که مایع از دست می‌دهد (اسهال / استفراغ)، ریسک هایپوتنشن بالاتر است \oudarr{} می‌توانیم مایع‌درمانی کنیم.
\item لوفکسیدین احتمال \en{BP} $\downarrow$ پایین‌تری دارد.
\end{itemize}

% pic14.jpg = 20261013_024415038.jpg; handwritten page 14.
\oudpage{14}{\rule{4cm}{0.2pt}}
\begin{itemize}
\item اضطراب و بی‌قراری:
\begin{itemize}
\item کلونیدین کمک می‌کند ولی کافی نیست.
\item \en{BNZ} می‌دهیم.
\begin{latin}
Clonazepam 0.5--1 mg 2--4/d
\end{latin}
\item در سرپایی \en{BNZ} نمی‌دهیم.
\item آنتی‌هیستامین می‌دهیم: هیدروکسی‌زین و دیفن‌هیدرامین.
\end{itemize}
\item بی‌خوابی:
\begin{itemize}
\item \en{Z-type} \oudarr{} زولپیدم ۵–۱۰ میلی‌گرم.
\item آنتی‌سایکوتیک ۲ \oudarr{} کوئتیاپین ۲۰–۵۰ میلی‌گرم.
\end{itemize}
\item درد اسکلتی‌عضلانی \oudarr{} \en{NSAID} (کتورولاک، ایبوپروفن).
\item ضداسپاسم‌ها \oudarr{} متوکلوپرامید، سایکلوبنزاپرین.
\item دیسترس \en{GI} \oudarr{} ضدتهوع‌ها (اندانسترون، پروکلرپرازین).
\item اسهال \oudarr{} لوپرامید، اکتروتاید (مولد شبه میو).
\end{itemize}

\ifnum\OUDStandalone=1
  \def\OUDFinish{\end{document}}
\else
  \let\OUDFinish\relax
\fi
\OUDFinish
